Zero-Space Projection (\ZSP) is a core mechanism of the 
\abnqss v3.0 physics-native computing architecture, whose digital 
twin has demonstrated projection linearity $R^2 = 1.0$ and a 
null-space dimension formula $K = N(N-1)/2$. This paper 
systematically evaluates \ZSP on three perception tasks: 
single-frame multi-object tracking (\MOT), multi-frame \MOT, and 
cooperative perception. Our findings are threefold: 
(i) on \textbf{single-frame \MOT}, \ZSP reaches $99.96\%$ of the 
accuracy of the Hungarian algorithm but is $20{,}000\times$ slower; 
(ii) on \textbf{multi-frame \MOT}, a paired $t$-test across 30 seeds 
shows \ZSP is $6.7\%$ worse than Hungarian on average 
($p = 0.0001$, win rate $0\%$); 
(iii) on \textbf{cooperative perception}, all three methods 
(centralized Hungarian, distributed Hungarian, and distributed \ZSP) 
achieve $100\%$ accuracy, but distributed Hungarian is fastest with 
zero communication overhead. 

We analyze the root cause: linear assignment problems possess a 
\emph{dual structure} that enables the Hungarian algorithm to achieve 
exact optimality in polynomial time, whereas \ZSP is an iterative 
approximation that cannot, by construction, outperform the exact 
algorithm. We conclude that the value of \ZSP---and of physics-native 
computing more broadly---lies not in \emph{replacing} problems already 
solved by exact algorithms, but in \emph{solving problems without good 
algorithms}: material simulation, physics-layer cryptography, 
high-precision sensing, and variational field solving. This work 
provides an explicit boundary for physics-native computing paradigms.